\section{The native row lattice and the simultaneous Bessel identity}

% BEGIN B87H1
\paragraph{B87H1. Finite generation of the actual sharp and base row lattices.}
Fix an original family, $x>\mathcal F.x_1$, and $i\in\mathcal F.I_x$.
Write $K=D.\mathrm{what}_2.\mathrm{LL}$, $\mathcal O=\mathcal O_K$,
and $j=D.j_{EE}:W\longrightarrow E$ for the coefficient embedding of the
sharp row space. For every tuple $t$ put
\[
 A_t=\mathcal O\langle y\in W:j(y)\in\mathrm{RRset}(t)\rangle,
 \qquad T=\{\mu_{\rm At}(w):w\text{ a finite place}\}.
\]
The set $T$ is finite: away from the finite support of this member
the tuple is $\mu_0$. The row module $A_t$ is finitely generated, even
before requiring its vectors to span $W$. Indeed let $P_t$ be orthogonal
projection onto $Q_t$ and put $C_t=\iota_tP_t$. The native definition
of $\mathrm{RRset}(t)$ uses $C_t$ on the intersection of the integral
carrier with $V_{0,t}$. The carrier is generated over the row
coefficient ring by its distinguished finite list $g_1,\dots,g_s$.
The image of that coefficient ring in $K$ is integral over
$\mathbf Z$, hence belongs to $\mathcal O$. Consequently
$j(A_t)$ is a submodule of $\sum_{a=1}^s\mathcal O C_tg_a$ in $E$,
where scalars act through the distinguished embedding. The latter is
a finite module over a Noetherian ring. The independence of the
images of the native rational basis makes $j$ injective. This proves
the assertion about $A_t$.

Put $A=\sum_{t\in T}A_t$, again finitely generated. The defining
family equality for $R=D.\mathrm{div}.\mathrm{RR}.L$ says that every
$z\in R$ lies in $(A_{\mu_{\rm At}(w)})_w$ at every $w$, hence in
$A_w$ at every $w$. To see $z\in A$, consider the ideal
$J=\{b\in\mathcal O:bz\in A\}$. At each finite prime $J$ contains
an element outside that prime. A proper ideal of $\mathcal O$
would be contained in a maximal ideal, and all maximal ideals here
are finite primes. Thus $J=\mathcal O$. It follows that $R\subset A$,
so $R$ is finitely generated. For the original base row
$\mathcal F.\mathrm{ex}(x,h,i).p_0$ the stored base row
definition gives the same localization criterion, with its own
$\iota_t,P_t,j$ and integral generators. It gives finite
generation there as well.
% END B87H1

% BEGIN B87H2
\paragraph{B87H2. The exact rank still required at the used tuples.}
With the same definitions on either side,
\[
 K R=\bigcap_{t\in T}K A_t\quad\text{inside }W.
\]
For $z\in R$ and each attained $t$ choose a place with that tuple.
Its local denominator is nonzero and sends $z$ into $A_t$.
Conversely, for $z$ in the intersection, clear the finitely many
coefficients in an expression for $z$ in each $K A_t$.
The product $b\ne0$ of the resulting denominators has
$bz\in A_t$ for every $t$, hence $bz\in R$ by its defining
equality, using local multiplier $1$ everywhere.

In particular sharp $\mathrm{RRFull}$ is now equivalent to the
following finite-rank question on this very member. For each
$t\in T$ find $d=\dim_K W$ vectors $y_a\in W$ with
$j(y_a)\in\mathrm{RRset}(t)$ and nonzero coordinate determinant
in the original rational basis. Indeed a generating set for a
spanning subspace of dimension $d$ contains a basis. The tuples
may use different bases. If $d=0$ the empty basis works. All these
facts use the recorded $j$ without changing its rational structure.
The determinant assertion for every tuple on the eventual family
range is the outstanding rank input in this argument.
% END B87H2

% BEGIN B87H3
\paragraph{B87H3. Detecting the one common Bessel isometry by finitely many scalar identities.}
Fix any native datum $D$. Choose orthonormal bases in $E=D.\mathrm{what}_2.EE$
and in each $M_u=D.\mathrm{what}_2.\mathrm{Mu}(u)$.
Write the two native pairings as $F^{\mathsf T}B_u\phi$ and
$F^{\mathsf T}S_u\phi$ for $\mathrm{Bu}$ and $\mathrm{besS}$,
respectively. Both are complex bilinear: the stored $\mathrm{Bu}$
is linear in each variable, and the conjugate-linearity of $s_{EE}$
cancels that of the first inner-product variable in $\mathrm{besS}$.
The map on its second variable is complex linear.
For each matrix put $G(B_u)=B_uB_u^*$ and likewise for $S_u$.
Normalize each nonzero such Gram matrix by its trace, writing
$H(B_u)=G(B_u)/\operatorname{tr}G(B_u)$. With $d=\dim_{\mathbf C} E>0$,
the unitary clause of the original Bessel identity is equivalent
to the following two tests:
\begin{enumerate}
\item $B_u=0$ exactly when $S_u=0$, for every $u$.
\item For all strings $u_1,\ldots,u_k$ of nonzero rows,
$1\le k\le 2d^2-1$,
\[
 \operatorname{tr}\big(H(B_{u_1})\cdots H(B_{u_k})\big)
 =\operatorname{tr}\big(H(S_{u_1})\cdots H(S_{u_k})\big).
\]
\end{enumerate}
The word set indexing the rows is finite. If $d=0$ the common
isometry and all row isometries can be identities and every pairing
vanishes.

To prove the test, the native equality in these coordinates reads
$B_u=a_u V^{\mathsf T}S_uU_u$, with $V$ unitary and
$a_u\ne0$. This implies both assertions by cancellation and
invariance of trace. For the converse, take the paired matrices
$H(B_u)\oplus H(S_u)$ acting on $\mathbf C^d\oplus\mathbf C^d$.
Their unital algebra is closed under adjoints because the
generators are Hermitian. It lies in $M_d(\mathbf C)\oplus M_d(\mathbf C)$,
of dimension $2d^2$. The spans of words of length at most $k$
start with dimension one. Equality of two consecutive spans makes
that span stable under right multiplication by every generator.
Hence words through length $2d^2-1$ span the whole algebra.
The two coordinate traces agree on this span, including the
identity. In the finite-dimensional complex star algebra each simple
matrix summand $M_r(\mathbf C)$ acts in a unitary representation as
copies of its standard $r$-dimensional module: its diagonal matrix
units split the space orthogonally, and the off-diagonal units
identify the pieces isometrically. The trace of its central
identity is $r$ times this multiplicity. Equality of the traces
therefore gives one unitary $W$ carrying the second coordinate
representation to the first.

For each nonzero row put
$a_u=(\operatorname{tr}G(B_u)/\operatorname{tr}G(S_u))^{1/2}$
and $C_u=a_u W S_u$. Then $C_uC_u^*=B_uB_u^*$.
The assignment $C_u^*v\mapsto B_u^*v$ preserves inner products
and extends by orthonormal bases to a unitary on $M_u$.
Taking its adjoint supplies $U_u$ with $B_u=C_u U_u$.
Take $V=W^{\mathsf T}$ once for all rows.
Zero rows allow $a_u=1$ and $U_u=1$.

The remaining rationality clause has a separate finite test. On the
two given $K$-bases $n_{ua},e_b$ check
\[
 \mathrm{Bu}_u(j_{N,u}n_{ua}, j_{EE} e_b)
   =\mathrm{emb}_{K}\bigl((D.\mathrm{div}.Bk)_u(n_{ua})(e_b)\bigr)
\]
for every $u,a,b$. Semilinearity of the two coordinate embeddings
and bilinearity of the native pairings extend it to every pair of
rational vectors. Combined with the zero and mixed-trace tests, these
are necessary and sufficient for $\mathrm{BesselId}(D)$.
Proving their values for the original periods at every sufficiently
large member remains part of the period and height obligation.
% END B87H3
